
% \iffalse meta-comment
%
% Copyright (C) 2014-2019 by Tibor Tomacs
%
% This file may be distributed and/or modified under the
% conditions of the LaTeX Project Public License, either version 1.2
% of this license or (at your option) any later version.
% The latest version of this license is in:
%
% http://www.latex-project.org/lppl.txt
%
% and version 1.2 or later is part of all distributions of LaTeX
% version 1999/12/01 or later.
%
% \fi
%
% \iffalse
%<*driver>
\ProvidesFile{thesis-ekf.dtx}
%</driver>
%<class>\NeedsTeXFormat{LaTeX2e}[1999/12/01]
%<class>\ProvidesClass{thesis-ekf}[2019/05/06 v3.1 Thesis class for Eszterhazy Karoly University (Eger, Hungary)]
%<class>
%<*driver>
\documentclass{ltxdoc}
\OnlyDescription
\usepackage[utf8]{inputenc}
\usepackage[T1]{fontenc}
\usepackage[a4paper,left=2in]{geometry}
\PassOptionsToPackage{active=onlycs}{magyar.ldf}
\usepackage[english,magyar]{babel}
\frenchspacing
\usepackage{paralist,textcomp,fancyvrb}
\setlength{\MacroTopsep}{0pt}
\fvset{gobble=1,fontsize=\footnotesize,commandchars=+<>}
\newcommand{\param}[1]{\hspace*{-9pt}{\small#1}}
\begin{document}
	\DocInput{./thesis-ekf.dtx}
\end{document}
%</driver>
% \fi
%
% \CharacterTable
% {Upper-case \A\B\C\D\E\F\G\H\I\J\K\L\M\N\O\P\Q\R\S\T\U\V\W\X\Y\Z
%     Lower-case \a\b\c\d\e\f\g\h\i\j\k\l\m\n\o\p\q\r\s\t\u\v\w\x\y\z
%     Digits \0\1\2\3\4\5\6\7\8\9
%     Exclamation         \!     Double quote    \"     Hash (number)   \#
%     Dollar              \$     Percent         \%     Ampersand       \&
%     Acute accent        \'     Left paren      \(     Right paren     \)
%     Asterisk            \*     Plus            \+     Comma           \,
%     Minus               \-     Point           \.     Solidus         \/
%     Colon               \:     Semicolon       \;     Less than       \<
%     Equals              \=     Greater than    \>     Question mark   \?
%     Commercial at       \@     Left bracket    \[     Backslash       \\
%     Right bracket       \]     Circumflex      \^     Underscore      \_
%     Grave accent        \`     Left brace      \{     Vertical bar    \|
%     Right brace         \}     Tilde           \~}
%
% \GetFileInfo{thesis-ekf.cls}
% \title{Thesis class for\\ Eszterházy Károly University\\\textsf{thesis-ekf.cls}\\[2mm]{\normalsize Version 3.1\\2019/05/06\\}}
% \author{Tibor Tómács\\{\small\ttfamily tomacs.tibor@uni-eszterhazy.hu}}
% \date{}
% \maketitle
%
% \selectlanguage{english}
% This is a class file for theses and dissertations at Eszterházy Károly University, Eger, Hungary (website: \texttt{https://uni-eszterhazy.hu}). The documentation is in Hungarian.
%
% \selectlanguage{magyar}
% \section{Bevezető}
%
% A \texttt{thesis-ekf} dokumentumosztály segítségével az Eszterházy Károly Egyetem szabályzatának megfelelő szakdolgozatokat lehet készíteni. A névben szereplő \texttt{ekf} az Eszterházy Károly Főiskola rövidítése, ugyanis ez a dokumentumosztály 2014-től létezik, amikor az intézmény még főiskola volt. A formai követelmények a következők:
%
% \medskip
% \begin{compactitem}
% \item A4-es lap- és 12\,pt betűméret,
% \item a margó a kötés oldalon 30\,mm, a többi 25\,mm,
% \item oldalszámozás a láblécben arab számozással,
% \item a fejezetcímek középre, a további szintek címei balra igazítva,
% \item a főszöveg antikva betűcsaláddal kiszedve,
% \item sorkizárt igazítás, másfeles sortávolság.
% \end{compactitem}
%
% \medskip\noindent
% A \texttt{thesis-ekf} ezeket a paramétereket automatikusan beállítja, továbbá a megfelelő címoldal elkészítését is elvégzi.
%
% \section{A dokumentumosztály betöltése}
%
% A \texttt{thesis-ekf} dokumentumosztály a \texttt{report} osztályt használja alapként, továbbá a következő csomagokat tölti be: \texttt{cmap}, \texttt{lmodern}, \texttt{fixcmex}, \texttt{kvoptions}, \texttt{etoolbox}, \texttt{ifpdf}, \texttt{setspace}, \texttt{hyperref}, \texttt{geometry}, \texttt{graphicx}, \texttt{upqute}.
%
% \medskip
% \begin{macro}{\documentclass}\param{\oarg{opciók}\texttt{\{thesis-ekf\}}}
% A dokumentumosztályt a preambulum elején ezzel a paranccsal lehet betölteni.
% \end{macro}
%
% \medskip\noindent
% Opciók nélkül akkor használja, ha a szakdolgozatot egyoldalasan szeretné kinyomtatni, vagy ha olyan elektronikus verziót akar, amelyben a linkek nem színes karakterrel jelennek meg. A fontosabb \meta{opciók} a következők:
%
% \medskip
% \begin{macro}{twoside}
% Ha a szakdolgozatot kétoldalasan szeretné kinyomtatni, akkor ezt az opciót alkalmazza! Ne használja egyoldalas nyomtatáshoz illetve elektronikus verzióhoz!
% \end{macro}
%
% \begin{macro}{colorlinks}
% A linkek színes karakterekkel jelennek meg. Ezt csak a szakdolgozat elektronikus verziójához használja, a nyomtatott verzióhoz nem kell!
% \end{macro}
%
% \begin{macro}{tocnopagenum}
% A tartalomjegyzéknek nem lesz oldalszámozása. Ha közvetlenül a címoldalt követően van elhelyezve a tartalomjegyzék, akkor az első számozott oldal csak ezután következik.
% \end{macro}
%
% \section{Címoldal létrehozása}
%
% \begin{macro}{\institute}\param{\marg{intézmény neve}}
% Ezzel kell megadni annak az intézménynek a nevét, ahol a szakdolgozat készült. Az Eszterházy Károly Egyetem esetében az egyetem nevét nem kell kiírni, mert azt tartalmazza a logója. Ekkor elég csak az intézet nevét feltüntetni.
% \end{macro}
%
% \begin{macro}{\title}\param{\marg{szakdolgozat címe}}
% Ezzel kell megadni a szakdolgozat címét.
% \end{macro}
%
% \begin{macro}{\author}\param{\texttt{\{}\meta{szerző neve}\texttt{\textbackslash\textbackslash}\meta{szak}\texttt{\}}}
% Ezzel kell megadni a szakdolgozat szerzőjének a nevét és a szakot.
% \end{macro}
%
% \begin{macro}{\supervisor}\param{\texttt{\{}\meta{témavezető neve}\texttt{\textbackslash\textbackslash}\meta{beosztás}\texttt{\}}}
% Ezzel kell megadni a szakdolgozat témavezetőjének a nevét és beosztását.
% \end{macro}
%
% \begin{macro}{\city}\param{\marg{város}}
% Ezzel kell megadni annak a városnak a nevét, ahol az intézmény található.
% \end{macro}
%
% \begin{macro}{\date}\param{\marg{dolgozat leadásának éve}}
% Ezzel kell megadni a szakdolgozat leadásának az évét. Ha nem adja meg, akkor az aktuális évszám fog megjelenni.
% \end{macro}
%
% \begin{macro}{\maketitle}
% A standard dokumentumosztályokhoz hasonlóan a címoldal itt is ezzel a paranccsal hozható létre.
% \end{macro}
%
% \section{Példa a használatra}
%
% A következő sablon akkor működik helyesen, ha UTF-8 kódolású fájlban van.
%
%\begin{Verbatim}
%\documentclass{thesis-ekf}
%\usepackage[utf8]{inputenc}% Ez a sor csak 2018 előtt telepített TeX-rendszer esetén kell.
%\usepackage[T1]{fontenc}
%\PassOptionsToPackage{defaults=hu-min}{magyar.ldf}
%\usepackage[magyar]{babel}
%
%\begin{document}
%
%\institute{Matematikai és Informatikai Intézet}
%\title{Szakdolgozat címe}
%\author{Hallgató neve\\szak}
%\supervisor{Konzulens neve\\beosztás}
%\city{Eger}
%\date{2019}
%\maketitle
%
%\tableofcontents
%\chapter{Fejezet címe}
%\section{Szakasz címe}
%\begin{thebibliography}{1}
%\bibitem{cimke} \textsc{Szerző}: Cím, Kiadó, Hely, évszám.
%\end{thebibliography}
%
%\end{document}
%\end{Verbatim}
%
% \section{A dokumentumosztály testreszabása}
%
% A \texttt{thesis-ekf} dokumentumosztály könnyen testreszabható opciókkal és parancsokkal, így más egyetemek is használhatják.
%
% \subsection{A dokumentumosztály opciói}
%
% A |twoside|, |colorlinks| és |tocnopagenum| opciókról már volt szó. További opciók:
%
% \medskip
% \begin{macro}{centeredchapter}\param{\texttt{=false}}
% Ezzel a fejezetcímek nem középre, hanem balra zártan lesznek kiszedve.
% \end{macro}
%
% \begin{macro}{warning}\param{\texttt{=false}}
% A dokumentumosztály figyelmeztetéseinek kikapcsolása.
% \end{macro}
%
% \begin{macro}{fontsize}\param{\texttt{=}\meta{betűméret}}
% A \meta{betűméret} lehetséges értékei |10pt|, |11pt| és |12pt| (alapérték |12pt|). Ez adja meg az alapbetűméretet.
% \end{macro}
%
% \begin{macro}{logodown}
% Alapbeállítás esetén a logó az intézmény neve felett van. Ezzel a logó az intézmény neve alá kerül.
% \end{macro}
%
% \begin{macro}{logofont}\param{\texttt{=}\meta{betűtípus}}
% Ha a logó helyén szöveg van -- lásd a |\logo| parancsot --, akkor annak a betűtípusa (alapérték |\large\scshape|).
% \end{macro}
%
% \begin{macro}{logosep}\param{\texttt{=}\meta{méret}}
% A logó alatti térköz a normál sortávolságon felül, ha a \texttt{logodown} opció nincs bekapcsolva (alapérték |0mm|).
% \end{macro}
%
% \begin{macro}{institutefont}\param{\texttt{=}\meta{betűtípus}}
% Az intézménynév betűtípusa (alapérték |\large\scshape|).
% \end{macro}
%
% \begin{macro}{institutesep}\param{\texttt{=}\meta{méret}}
%  Az intézmény neve alatti térköz a normál sortávolságon felül, ha a \texttt{logodown} opció be van kapcsolva (alapérték |10mm|).
% \end{macro}
%
% \begin{macro}{titlefont}\param{\texttt{=}\meta{betűtípus}}
% A dolgozat címének betűtípusa (alapérték |\Huge\bfseries|).
% \end{macro}
%
% \begin{macro}{titlesep}\param{\texttt{=}\meta{méret}}
% A cím feletti térköz mérete |\stretch{1}|. Ezzel a cím alatti térközt lehet beállítani a normál sortávolságon felül (alapérték |\stretch{1.5}|).
% \end{macro}
%
% \begin{macro}{captionfont}\param{\texttt{=}\meta{betűtípus}}
% A szerző és témavezető nevei feletti feliratok betűtípusa (alapérték |\large\bfseries|).
% \end{macro}
%
% \begin{macro}{captionsep}\param{\texttt{=}\meta{méret}}
% Az előbbi feliratok alatti térköz a normál sortávolságon felül (alapérték |0mm|).
% \end{macro}
%
% \begin{macro}{authorfont}\param{\texttt{=}\meta{betűtípus}}
% A szerző nevének és szakjának, illetve a témavezető nevének és beosztásának betűtípusa (alapérték |\large|).
% \end{macro}
%
% \begin{macro}{authoralign}\param{\texttt{=}\meta{igazítás}}
% A szerző adatait tartalmazó doboz tartalmának igazítása. Az \meta{igazítás} lehetséges értékei |left| (balra igazítás, ez az alapérték), |center| (középre igazítás) |right| (jobbra igazítás).
% \end{macro}
%
% \begin{macro}{supervisoralign}\param{\texttt{=}\meta{igazítás}}
% A témavezető adatait tartalmazó doboz tartalmának igazítása. Az \meta{igazítás} lehetséges értékei |left| (balra igazítás, ez az alapérték), |center| (középre igazítás) |right| (jobbra igazítás).
% \end{macro}
%
% \begin{macro}{authorxmargin}\param{\texttt{=}\meta{méret}}
% Ezzel a szerzőnél és a témavezetőnél egy extra margóméret adható meg, azaz a normál margó ennyivel nő (alapérték |10mm|).
% \end{macro}
%
% \begin{macro}{authorsep}\param{\texttt{=}\meta{méret}}
% A szerző illetve témavezető alatti térköz a normál sortávolságon felül (alapérték |15mm|).
% \end{macro}
%
% \begin{macro}{cityfont}\param{\texttt{=}\meta{betűtípus}}
% A város nevének betűtípusa (alapérték |\large\scshape|).
% \end{macro}
%
% \begin{macro}{datefont}\param{\texttt{=}\meta{betűtípus}}
% Az évszám betűtípusa (alapérték |\large\scshape|).
% \end{macro}
%
% \begin{macro}{datesep}\param{\texttt{=}\meta{elválasztás}}
% A város és az évszám közötti elválasztás (alapérték |{,~}|). Például |datesep=\\| esetén a városnév alá kerül az évszám.
% \end{macro}
%
% \subsection{Parancsok}
%
% \begin{macro}{\setkeys}\param{\texttt{\{thesis-ekf\}}\marg{opciók}}
% A \texttt{thesis-ekf} opciói a |fontsize| kivételével, ezzel is beállíthatók. Például |\setkeys{thesis-ekf}{logodown,tocnopagenum}|. Ha a |\setkeys| parancsot a konfigurációs fájlba írja (lásd később), akkor abba a |fontsize| opció is beírható. Azok az opciók, melyekben parancs van, mint például a betűtípusra vonatkozók, csak a |\setkeys| parancsba írhatók!
% \end{macro}
%
% \begin{macro}{\hypersetup}\param{\marg{hyperref opciók}}
% A \texttt{hyperref} csomag opciói ezzel állíthatók be. Például, ha a linkek színét pirosra akarja állítani: |\hypersetup{allcolors=red}|.
% \end{macro}
%
% \begin{macro}{\geometry}\param{\marg{geometry opciók}}
% A \texttt{geometry} csomag opciói ezzel állíthatók be. Például, ha B5 lapméretet szeretne, 20\,mm margókkal, kivéve a belső margót, amely 25\,mm:\\ |\geometry{b5paper,top=20mm,bottom=20mm,inner=25mm,outer=20mm}|.
% \end{macro}
%
% \begin{macro}{\singlespacing}
% Alapbeállítás esetén a sortávolság másfeles. A \texttt{setspace} csomag |\singlespacing| parancsával visszaállítható a normál méretű sortávolság.
% \end{macro}
%
% \begin{macro}{\logo}\param{\marg{kép vagy szöveg}}
% Ezt a parancsot a |\maketitle| előtt vagy a konfigurációs fájlban (lásd később) kell használni. Ezzel kell megadni az intézmény logóját. Például |\logo{\includegraphics[width=9cm]{logo}}|. Ha nem adja meg, akkor az Eszterházy Károly Egyetem angol, német vagy magyar logója fog megjelenni aszerint, hogy a \texttt{babel} csomaggal melyik nyelvet töltötte be. Ha nem akar logót, akkor írja be a |\logo{}| parancsot.
% \end{macro}
%
% \begin{macro}{\authorcaption}\param{\marg{szerző neve feletti felirat}}
% Ezt a parancsot a |\maketitle| előtt vagy a konfigurációs fájlban (lásd később) kell használni. Ezzel adhatja meg a szakdolgozat szerzőjének neve feletti feliratot. Alapértéke \texttt{Author}, magyar nyelv esetén \texttt{Készítette}, német nyelv esetén \texttt{Autor}.
% \end{macro}
%
% \begin{macro}{\supervisorcaption}\param{\marg{témavezető neve feletti felirat}}
% Ezt a parancsot a |\maketitle| előtt vagy a konfigurációs fájlban (lásd később) kell használni. Ezzel adhatja meg a szakdolgozat témavezetőjének neve feletti feliratot. Alapértéke \texttt{Supervisor}, magyar nyelv esetén \texttt{Témavezető}, német nyelv esetén \texttt{Betreuer}.
% \end{macro}
%
% \subsection{Konfigurációs fájl}
%
% \begin{macro}{thesis-ekf.cfg}
% Az átparaméterező opciókat és parancsokat célszerű egy \texttt{thesis-ekf.cfg} fájlba írni, amit a \texttt{tex} forrásfájl mellé kell tenni. A dokumentumosztály betöltésekor ez a fájl automatikusan betöltődik. A konfigurációs fájlban az ékezetes betűket repülő ékezettel érdemes beírni, hogy minden kódolású \texttt{tex} fájl esetén működjön.
% \end{macro}
%
% \begin{macro}{\AtEndOfClass}\param{\marg{parancsok}}
% A konfigurációs fájlban minden parancsot a |\setkeys| kivételével ebbe a parancsba kell írni.
% \end{macro}
%
% \subsection{Egy példa az átparaméterezésre}
%
% A következőket írja be a \texttt{thesis-ekf.cfg} fájlba, majd tegye a \texttt{tex} forrásfájl mellé.
%
%\begin{Verbatim}
%\setkeys{thesis-ekf}{fontsize=11pt,centeredchapter=false,titlefont=\Huge\sffamily}
%\AtEndOfClass{
%  \hypersetup{allcolors=red}
%  \geometry{b5paper,top=20mm,bottom=20mm,inner=25mm,outer=20mm}
%  \singlespacing
%  \authorcaption{\+textquotesingle{I}rta}
%}
%\end{Verbatim}
%
% \StopEventually{}
%    \begin{macrocode}
\@ifundefined{pdffontattr}{}{\ifnum\pdfoutput>0\RequirePackage{cmap}\fi}
\RequirePackage{lmodern,fixcmex,kvoptions,etoolbox,ifpdf,setspace,graphicx}

\SetupKeyvalOptions{family=thesis-ekf,prefix=thesisekf@}
\DeclareVoidOption{twoside}{\PassOptionsToClass{twoside}{report}}
\DeclareVoidOption{colorlinks}{\PassOptionsToPackage{colorlinks}{hyperref}}
\DeclareBoolOption{logodown}
\DeclareBoolOption{tocnopagenum}
\DeclareBoolOption[true]{centeredchapter}
\DeclareBoolOption[true]{warning}
\DeclareStringOption[12pt]{fontsize}
\DeclareStringOption[10mm]{institutesep}
\DeclareStringOption[0mm]{logosep}
\DeclareStringOption[\stretch{1.5}]{titlesep}
\DeclareStringOption[15mm]{authorsep}
\DeclareStringOption[10mm]{authorxmargin}
\DeclareStringOption[0mm]{captionsep}
\DeclareStringOption[left]{authoralign}
\DeclareStringOption[left]{supervisoralign}
\DeclareStringOption[\large\scshape]{institutefont}
\DeclareStringOption[\large\scshape]{logofont}
\DeclareStringOption[\Huge\bfseries]{titlefont}
\DeclareStringOption[\large]{authorfont}
\DeclareStringOption[\large\bfseries]{captionfont}
\DeclareStringOption[\large\scshape]{cityfont}
\DeclareStringOption[\large\scshape]{datefont}
\DeclareStringOption[,~]{datesep}
\ProcessKeyvalOptions{thesis-ekf}

\InputIfFileExists{thesis-ekf.cfg}{}{}

\def\thesisekf@ten{10pt}
\def\thesisekf@eleven{11pt}
\def\thesisekf@twelve{12pt}

\ifx\thesisekf@fontsize\thesisekf@ten\else
\ifx\thesisekf@fontsize\thesisekf@eleven\else
\ifx\thesisekf@fontsize\thesisekf@twelve\else
\ClassWarning{thesis-ekf}{Unused option: fontsize=\thesisekf@fontsize. Use fontsize=10pt, fontsize=11pt or fontsize=12pt!}
\def\thesisekf@fontsize{12pt}
\fi\fi\fi

\PassOptionsToClass{\thesisekf@fontsize}{report}
\LoadClass{report}

\RequirePackage[unicode,linktocpage,allcolors=blue,pdfstartview=FitH,bookmarksnumbered,pdfborder={0 0 0}]{hyperref}
\RequirePackage[a4paper,top=25mm,bottom=25mm,inner=30mm,outer=25mm]{geometry}

\onehalfspacing

\AtEndPreamble{\RequirePackage{upquote}}

\AfterEndPreamble{

\ifthesisekf@tocnopagenum
  \hypersetup{pageanchor=false}
  \let\thesisekf@tableofcontents\tableofcontents
  \def\tableofcontents{%
     \global\let\thesisekf@thepage\thepage%
     \global\let\thepage\empty%
     \thesisekf@tableofcontents%
     \clearpage%
     \global\c@page\@ne%
     \global\let\thepage\thesisekf@thepage%
     \hypersetup{pageanchor}}
\fi

\ifthesisekf@centeredchapter
  \let\thesisekf@makechapterhead\@makechapterhead
  \let\thesisekf@makeschapterhead\@makeschapterhead
  \let\thesisekf@raggedright\raggedright
  \def\@makechapterhead#1{
    \let\raggedright\centering
    \thesisekf@makechapterhead{#1}
    \let\raggedright\thesisekf@raggedright}
  \def\@makeschapterhead#1{
    \let\raggedright\centering
    \thesisekf@makeschapterhead{#1}
    \let\raggedright\thesisekf@raggedright}
\fi

\def\@dottedtocline#1#2#3#4#5{%
  \ifnum #1>\c@tocdepth \else
    \vskip \z@ \@plus.2\p@
    {\leftskip #2\relax \rightskip \@tocrmarg \parfillskip -\rightskip
     \parindent #2\relax\@afterindenttrue
     \interlinepenalty\@M
     \leavevmode
     \@tempdima #3\relax
     \advance\leftskip \@tempdima \null\nobreak\hskip -\leftskip
     {#4}\nobreak
     \leaders\hbox{$\m@th
        \mkern \@dotsep mu\hbox{.}\mkern \@dotsep
        mu$}\hfill
     \nobreak
     \hb@xt@\@pnumwidth{\hfil\normalfont \normalcolor #5}%
     \par}%
  \fi}

\renewcommand*\l@chapter[2]{%
  \ifnum \c@tocdepth >\z@
    \addpenalty\@secpenalty
    \addvspace{1.0em \@plus\p@}%
    \setlength\@tempdima{1.8em}%
    \begingroup
      \parindent \z@ \rightskip \@pnumwidth
      \parfillskip -\@pnumwidth
      \leavevmode \bfseries
      \advance\leftskip\@tempdima
      \hskip -\leftskip
      #1\nobreak\hfil \nobreak\hb@xt@\@pnumwidth{\hss #2}\par
    \endgroup
  \fi}
\renewcommand*\l@section{\@dottedtocline{1}{1.8em}{2.5em}}
\renewcommand*\l@subsection{\@dottedtocline{2}{4.3em}{3.2em}}

\if@thesisekf@nochanged@authorcaption@%
    \@ifundefined{l@magyar}{}{%
    \iflanguage{magyar}{\gdef\thesisekf@authorcaption{K\'{e}sz\'{\i}tette}}{}}%
    \@ifundefined{l@ngerman}{}{%
    \iflanguage{ngerman}{\gdef\thesisekf@authorcaption{Autor}}{}}%
    \@ifundefined{l@german}{}{%
    \iflanguage{german}{\gdef\thesisekf@authorcaption{Autor}}{}}\fi

\if@thesisekf@nochanged@supervisorcaption@%
    \@ifundefined{l@magyar}{}{%
    \iflanguage{magyar}{\gdef\thesisekf@supervisorcaption{T\'{e}mavezet\H{o}}}{}}%
    \@ifundefined{l@ngerman}{}{%
    \iflanguage{ngerman}{\gdef\thesisekf@supervisorcaption{Betreuer}}{}}%
    \@ifundefined{l@german}{}{%
    \iflanguage{german}{\gdef\thesisekf@supervisorcaption{Betreuer}}{}}\fi

\if@thesisekf@nochanged@logo@%
    \@ifundefined{l@magyar}{}{%
    \iflanguage{magyar}{\if@thesisekf@exists@logo@hu@\gdef\thesisekf@logo{\includegraphics{eszterhazy-logo-hu}}\fi}{}}%
    \@ifundefined{l@ngerman}{}{%
    \iflanguage{ngerman}{\if@thesisekf@exists@logo@de@\gdef\thesisekf@logo{\includegraphics{eszterhazy-logo-de}}\fi}{}}%
    \@ifundefined{l@german}{}{%
    \iflanguage{german}{\if@thesisekf@exists@logo@de@\gdef\thesisekf@logo{\includegraphics{eszterhazy-logo-de}}\fi}{}}\fi

}

\newif\if@thesisekf@nochanged@logo@\@thesisekf@nochanged@logo@true
\def\logo#1{\@thesisekf@nochanged@logo@false\gdef\thesisekf@logo{#1}}
\def\thesisekf@logo{}
\newif\if@thesisekf@exists@logo@en@
\ifpdf\IfFileExists{eszterhazy-logo-en.pdf}{\@thesisekf@exists@logo@en@true}{}
    \else\IfFileExists{eszterhazy-logo-en.eps}{\@thesisekf@exists@logo@en@true}{}\fi
\newif\if@thesisekf@exists@logo@hu@
\ifpdf\IfFileExists{eszterhazy-logo-hu.pdf}{\@thesisekf@exists@logo@hu@true}{}
    \else\IfFileExists{eszterhazy-logo-hu.eps}{\@thesisekf@exists@logo@hu@true}{}\fi
\newif\if@thesisekf@exists@logo@de@
\ifpdf\IfFileExists{eszterhazy-logo-de.pdf}{\@thesisekf@exists@logo@de@true}{}
    \else\IfFileExists{eszterhazy-logo-de.eps}{\@thesisekf@exists@logo@de@true}{}\fi
\if@thesisekf@exists@logo@en@\def\thesisekf@logo{\includegraphics{eszterhazy-logo-en}}\fi
\def\thesisekf@logo@{\thesisekf@logo%
    \ifx\thesisekf@logo\@empty\ifthesisekf@warning\ClassWarning{thesis-ekf}{There isn't logo!}\fi\fi}

\def\institute#1{\gdef\thesisekf@institute{#1}}
\institute{}
\def\thesisekf@institute@{\thesisekf@institute%
    \ifx\thesisekf@institute\@empty\ifthesisekf@warning\ClassWarning{thesis-ekf}{There isn't institute!}\fi\fi}

\def\title#1{\gdef\thesisekf@title{#1}}
\title{}
\def\thesisekf@title@{\thesisekf@title%
    \ifx\thesisekf@title\@empty\ifthesisekf@warning\ClassWarning{thesis-ekf}{There isn't title!}\fi\fi}

\def\author#1{\gdef\thesisekf@author{#1}}
\author{}
\def\thesisekf@author@{\thesisekf@author%
    \ifx\thesisekf@author\@empty\ifthesisekf@warning\ClassWarning{thesis-ekf}{There isn't author!}\fi\fi}

\newif\if@thesisekf@nochanged@authorcaption@\@thesisekf@nochanged@authorcaption@true
\def\authorcaption#1{\@thesisekf@nochanged@authorcaption@false\gdef\thesisekf@authorcaption{#1}}
\def\thesisekf@authorcaption{Author}
\def\thesisekf@authorcaption@{%
    \ifx\thesisekf@author\@empty\else
    \ifx\thesisekf@authorcaption\@empty\ifthesisekf@warning\ClassWarning{thesis-ekf}{There isn't authorcaption!}\fi
    \else\thesisekf@authorcaption\fi\fi}

\def\supervisor#1{\gdef\thesisekf@supervisor{#1}}
\supervisor{}
\def\thesisekf@supervisor@{\thesisekf@supervisor%
    \ifx\thesisekf@supervisor\@empty\ifthesisekf@warning\ClassWarning{thesis-ekf}{There isn't supervisor!}\fi\fi}

\newif\if@thesisekf@nochanged@supervisorcaption@\@thesisekf@nochanged@supervisorcaption@true
\def\supervisorcaption#1{\@thesisekf@nochanged@supervisorcaption@false\gdef\thesisekf@supervisorcaption{#1}}
\def\thesisekf@supervisorcaption{Supervisor}
\def\thesisekf@supervisorcaption@{%
    \ifx\thesisekf@supervisor\@empty\else
    \ifx\thesisekf@supervisorcaption\@empty\ifthesisekf@warning\ClassWarning{thesis-ekf}{There isn't supervisorcaption!}\fi
    \else\thesisekf@supervisorcaption\fi\fi}

\def\city#1{\gdef\thesisekf@city{#1}}
\city{}
\def\thesisekf@city@{\thesisekf@city%
    \ifx\thesisekf@city\@empty\ifthesisekf@warning\ClassWarning{thesis-ekf}{There isn't city!}\fi\fi}

\def\date#1{\gdef\thesisekf@date{#1}}
\date{\number\year}
\def\thesisekf@date@{\thesisekf@date%
    \ifx\thesisekf@date\@empty\ifthesisekf@warning\ClassWarning{thesis-ekf}{There isn't date!}\fi\fi}

\def\thesisekf@datesep@{\ifx\thesisekf@city\@empty\else\ifx\thesisekf@date\@empty\else\thesisekf@datesep\fi\fi}

\newif\if@thesisekf@noempty@author@supervisor@\@thesisekf@noempty@author@supervisor@true

\def\maketitle{
\thispagestyle{empty}
{\centering
\ifthesisekf@logodown
    {\normalfont\normalsize\thesisekf@institutefont\thesisekf@institute@%
    \ifx\thesisekf@institute\@empty\else\\
        \ifx\thesisekf@logo\@empty\else\vspace{\thesisekf@institutesep}\fi\fi}
    {\normalfont\normalsize\thesisekf@logofont\thesisekf@logo@%
    \ifx\thesisekf@logo\@empty\else\\\fi}
\else
    {\normalfont\normalsize\thesisekf@logofont\thesisekf@logo@%
    \ifx\thesisekf@logo\@empty\else\\
        \ifx\thesisekf@institute\@empty\else\vspace{\thesisekf@logosep}\fi\fi}
    {\normalfont\normalsize\thesisekf@institutefont\thesisekf@institute@%
    \ifx\thesisekf@institute\@empty\else\\\fi}
\fi
\vspace*{\fill}
{\normalfont\normalsize\thesisekf@titlefont\thesisekf@title@%
\ifx\thesisekf@title\@empty\else\\\vspace{\thesisekf@titlesep}\fi}
\ifx\thesisekf@supervisor\@empty\ifx\thesisekf@author\@empty\@thesisekf@noempty@author@supervisor@false\fi\fi
\if@thesisekf@noempty@author@supervisor@
{\normalfont\normalsize\thesisekf@authorfont
\def\thesisekf@align@center{center}%
\def\thesisekf@align@left{left}%
\def\thesisekf@align@right{right}%
\ifx\thesisekf@authoralign\thesisekf@align@center\begin{tabular}[t]{@{\hspace{\thesisekf@authorxmargin}}c@{}}\fi
\ifx\thesisekf@authoralign\thesisekf@align@left\begin{tabular}[t]{@{\hspace{\thesisekf@authorxmargin}}l@{}}\fi
\ifx\thesisekf@authoralign\thesisekf@align@right\begin{tabular}[t]{@{\hspace{\thesisekf@authorxmargin}}r@{}}\fi
{\normalfont\normalsize\thesisekf@captionfont\thesisekf@authorcaption@}\\[\thesisekf@captionsep]
\thesisekf@author@
\end{tabular}
\hfill
\ifx\thesisekf@supervisoralign\thesisekf@align@center\begin{tabular}[t]{@{}c@{\hspace{\thesisekf@authorxmargin}}}\fi
\ifx\thesisekf@supervisoralign\thesisekf@align@left\begin{tabular}[t]{@{}l@{\hspace{\thesisekf@authorxmargin}}}\fi
\ifx\thesisekf@supervisoralign\thesisekf@align@right\begin{tabular}[t]{@{}r@{\hspace{\thesisekf@authorxmargin}}}\fi
{\normalfont\normalsize\thesisekf@captionfont\thesisekf@supervisorcaption@}\\[\thesisekf@captionsep]
\thesisekf@supervisor@
\end{tabular}
\\\vspace{\thesisekf@authorsep}}
\fi
{\normalfont\normalsize\thesisekf@cityfont\thesisekf@city@\thesisekf@datesep@}%
{\normalfont\normalsize\thesisekf@datefont\thesisekf@date@}
\par
}
\newpage
\if@twoside\thispagestyle{empty}\hbox{}\newpage\fi}

%    \end{macrocode}
% \Finale
\endinput